% Artifact appendix, in the CGO artifact evaluation layout
% (https://ctuning.org/ae/submission_extra.html).
% Use: \appendix % Artifact appendix, in the CGO artifact evaluation layout
% (https://ctuning.org/ae/submission_extra.html).
% Use: \appendix % Artifact appendix, in the CGO artifact evaluation layout
% (https://ctuning.org/ae/submission_extra.html).
% Use: \appendix % Artifact appendix, in the CGO artifact evaluation layout
% (https://ctuning.org/ae/submission_extra.html).
% Use: \appendix \input{artifact_appendix} before the bibliography (CGO
% places the artifact appendix before the References, up to two pages).
% After archiving the artifact, set its DOI here:
\newcommand{\artifactdoi}{\textit{DOI to be assigned on archiving}}
\newcommand{\artifactrepo}{\url{https://github.com/wraith1995/Halide/tree/inductive_artifact}}
% Figure references that fall back to text when built outside the paper.
\newcommand{\artfig}[2]{\ifcsname r@#1\endcsname Figure~\ref{#1}\else #2\fi}
\newcommand{\artsec}[2]{\ifcsname r@#1\endcsname Section~\ref{#1}\else #2\fi}

\section{Artifact Appendix}
\sloppy

\subsection{Abstract}

The artifact is a fork of Halide with inductive functions, and the benchmark
suite behind \artfig{fig:cpuspeedup}{the CPU speedup figure},
\artfig{fig:gpuspeedup}{the GPU speedup figure} and
\artfig{fig:memuse}{the memory figure}. Each benchmark computes the same
recurrence twice: as an inductive function, and as a non-inductive
update definition over an RDom, which materializes the whole trajectory.
Each run also times the fastest third-party implementation that is
installed. One script regenerates the speedups and peak intermediate memory
of every app. A second runs an ablation of the CPU apps that adds the
optimizations inductive functions enable one at a time. Every form is checked
against a reference implementation before it is timed.

\subsection{Artifact check-list (meta-information)}

{\small
\begin{itemize}
  \item {\bf Algorithm:} inductive functions in Halide, scheduled with fusion,
    sliding window and storage folding.
  \item {\bf Program:} six CPU apps (RNG, biquads, sequence alignment, prefix
    mean, Viterbi, ODE) and two GPU apps (FlashAttention, Mamba-2 forward
    and backward), with their third-party baselines.
  \item {\bf Compilation:} Halide built against LLVM~22, which \texttt{uv}
    installs; C++17 at \texttt{-O3 -march=native}; CUDA~13 for the GPU apps.
  \item {\bf Data set:} generated by each app from a fixed seed.
  \item {\bf Run-time environment:} Linux x86-64 with jemalloc. The optional
    baselines are Boost.odeint, oneTBB, Intel IPP, parasail, Julia, scipy,
    PyTorch, Triton and mamba\_ssm.
  \item {\bf Hardware:} a multicore x86-64 CPU. The paper's machine
    (\artsec{sec:benchmarks}{in the paper}) has 32 cores with AVX-512; AVX2 also
    works. The GPU apps need an NVIDIA GPU with tensor cores.
  \item {\bf Execution:} two shell scripts. All-cores rows use one thread
    per physical core.
  \item {\bf Metrics:} best of 30 timed runs after 3 warmups (ms); peak
    intermediate heap (MB) from the Halide profiler.
  \item {\bf Output:} Markdown and text tables, plus logs.
  \item {\bf Experiments:} \artfig{fig:cpuspeedup}{CPU speedups},
    \artfig{fig:gpuspeedup}{GPU speedups}, \artfig{fig:memuse}{memory}, and
    the CPU ablation.
  \item {\bf How much disk space required (approximately)?:} 10~GB.
  \item {\bf How much time is needed to prepare workflow (approximately)?:}
    about an hour.
  \item {\bf How much time is needed to complete experiments (approximately)?:}
    a few hours; each app can also be run on its own.
  \item {\bf Publicly available?:} yes.
  \item {\bf Code licenses (if publicly available)?:} MIT (Halide). Vendored
    baselines keep their own licenses.
  \item {\bf Archived (provide DOI)?:} \artifactdoi.
\end{itemize}
}

\subsection{Description}

\subsubsection{How delivered}
The branch \texttt{inductive\_artifact} of \artifactrepo, archived at
\artifactdoi. The benchmark apps are in \path{apps/}, and the artifact's
scripts and instructions are in \path{inductive_artifact/}.

\subsubsection{Hardware dependencies}
The CPU apps run on any x86-64 Linux machine. The all-cores rows are
memory-bound, so their speedups depend on the memory bandwidth available per
core. The GPU apps need an NVIDIA GPU with tensor cores and CUDA~13.

\subsubsection{Software dependencies}
Linux, a C++17 compiler, GNU make, pkg-config, \texttt{uv}, jemalloc, and the
libpng, libjpeg and zlib headers. Every process the scripts launch, baselines
included, preloads jemalloc configured to keep freed memory, as in the paper.
The baselines are optional. Boost.odeint, oneTBB, Intel IPP and parasail are
found through \texttt{BOOST\_DIR}, \texttt{TBB\_DIR}, \texttt{IPP\_DIR} and
\texttt{PARASAIL\_DIR}, or a system install. Julia is found on the
\texttt{PATH}. The Python baselines come from
\path{inductive_artifact/requirements.txt}. When a baseline is missing, the
row reports the fastest baseline that is present.

\subsubsection{Data sets}
None; every input is generated.

\subsection{Installation}

{\footnotesize
\begin{verbatim}
git clone -b inductive_artifact \
  https://github.com/wraith1995/Halide.git
cd Halide
sudo apt install build-essential pkg-config \
  libpng-dev libjpeg-turbo8-dev zlib1g-dev \
  libjemalloc2 libboost-dev libtbb-dev
inductive_artifact/setup.sh
# or, with the Python and GPU baselines:
WITH_PYTHON_BASELINES=1 inductive_artifact/setup.sh
\end{verbatim}
}

\noindent\texttt{setup.sh} installs LLVM into \path{.venv} and builds the
Halide distribution, \path{distrib/}, which every app links against.

\subsection{Experiment workflow}

{\footnotesize
\begin{verbatim}
inductive_artifact/run_benchmarks.sh  # figures
inductive_artifact/run_ablation.sh    # ablation
\end{verbatim}
}

\noindent Both scripts accept \texttt{--only} with a comma-separated list of
apps: \texttt{cpu\_rng}, \texttt{cpu\_biquads}, \texttt{cpu\_alignment},
\texttt{prefixsum}, \texttt{viterbi}, \texttt{ode}, \texttt{flash\_attention}
and \texttt{cuda\_mamba2}. Finished rows are cached, so a partial run
updates the table rather than replacing it. The GPU apps print their own
ablation levels:

{\footnotesize
\begin{verbatim}
make -C apps/cuda_mamba2 SCAN=inductive \
  WMMA=true ABLATE_FUSION=1 test test_bwd
make -C apps/cuda_attention test
\end{verbatim}
}

\subsection{Evaluation and expected results}

\path{run_benchmarks.sh} writes \path{inductive_artifact/results/benchmarks.md},
with one row per app and configuration: the inductive and non-inductive
times, each form's peak intermediate memory, and the fastest baseline. The
all-cores CPU rows and the GPU rows are the configurations in the paper. The
single-thread rows and the Chebyshev row are controls that the paper does
not plot. The tables from the paper's machine are in
\path{apps/inductive_benchmarks_results.md} and
\path{apps/inductive_ablation_results/}. The claims to check:

\begin{enumerate}
  \item \textbf{Correctness.} Every form passes its app's check against the
    reference before it is timed.
  \item \textbf{Memory (\artfig{fig:memuse}{memory figure}).} The inductive
    forms allocate no intermediate memory, except for the backpointers that
    alignment (1094~MB) and Viterbi (134~MB) keep for their traceback. The
    non-inductive forms materialize the trajectory: 4295~MB for RNG,
    2147~MB for biquads, 10758~MB for alignment, 1074~MB each for prefix
    mean and ODE, and 671~MB for Viterbi. On the GPU, Mamba-2 forward needs
    42~MB against 141~MB. These values depend only on the problem sizes, so
    they should match exactly. The ODE and Viterbi all-cores rows run one
    batch element per core, so their memory scales with the core count;
    \texttt{HB\_NCORES=32} gives the paper's sizes.
  \item \textbf{CPU speedups (\artfig{fig:cpuspeedup}{CPU figure}).} The
    inductive form is faster than the non-inductive form in every app. On
    the paper's machine the speedups are prefix mean 2.00$\times$, Viterbi
    4.37$\times$, RNG 5.88$\times$, biquads 7.43$\times$, alignment
    16.47$\times$ and ODE 19.01$\times$. Against the baselines, the inductive
    form matches Julia (RNG) and oneTBB (prefix mean), reaches 93\% of
    Finding Fast Filters (biquads), and beats parasail (alignment) and
    Boost.odeint (ODE).
  \item \textbf{GPU speedups (\artfig{fig:gpuspeedup}{GPU figure}).} The
    inductive form beats the non-inductive form by 1.09$\times$
    (FlashAttention), 1.15$\times$ (Mamba-2 forward) and 1.04$\times$
    (Mamba-2 backward). Relative to the vendor kernels it reaches 94\% of
    PyTorch, 98\% of Triton on the forward pass and 73\% of Triton on the
    backward pass.
  \item \textbf{Ablation (CPU).} The levels are non-inductive, fusion only,
    fusion with a sliding window, and the full schedule with storage
    folding. Fusion alone is far slower than any other level, because each
    output recomputes the recurrence from its start. It is measured at a
    small size, where it runs 62--6038$\times$ slower than the full schedule.
    Adding the sliding window lands near the non-inductive form
    (1.3--14$\times$ the full schedule). ODE is the exception at
    39$\times$: its unfolded window is three times the non-inductive
    storage. Storage folding is the fastest level in every app.
\end{enumerate}

Absolute times depend on the machine. The all-cores speedups come from
memory traffic, so they are smaller on a machine with more memory bandwidth
per core. The ranking of the apps and the memory values should not change.

\subsection{Experiment customization}

Each app's Makefile takes its problem size on the command line, and the
README in each app lists the knobs. Environment variables set the number of
cores (\texttt{HB\_NCORES}), the core that single-thread rows are pinned to
(\texttt{HB\_PIN\_CORE}), the Python used for the baselines
(\texttt{HB\_PYTHON}), and the jemalloc library (\texttt{HB\_JEMALLOC}).

\subsection{Methodology}

Submission, reviewing and badging methodology:
\begin{itemize}
  \item \url{https://www.acm.org/publications/policies/artifact-review-and-badging-current}
  \item \url{https://cTuning.org/ae}
\end{itemize}
 before the bibliography (CGO
% places the artifact appendix before the References, up to two pages).
% After archiving the artifact, set its DOI here:
\newcommand{\artifactdoi}{\textit{DOI to be assigned on archiving}}
\newcommand{\artifactrepo}{\url{https://github.com/wraith1995/Halide/tree/inductive_artifact}}
% Figure references that fall back to text when built outside the paper.
\newcommand{\artfig}[2]{\ifcsname r@#1\endcsname Figure~\ref{#1}\else #2\fi}
\newcommand{\artsec}[2]{\ifcsname r@#1\endcsname Section~\ref{#1}\else #2\fi}

\section{Artifact Appendix}
\sloppy

\subsection{Abstract}

The artifact is a fork of Halide with inductive functions, and the benchmark
suite behind \artfig{fig:cpuspeedup}{the CPU speedup figure},
\artfig{fig:gpuspeedup}{the GPU speedup figure} and
\artfig{fig:memuse}{the memory figure}. Each benchmark computes the same
recurrence twice: as an inductive function, and as a non-inductive
update definition over an RDom, which materializes the whole trajectory.
Each run also times the fastest third-party implementation that is
installed. One script regenerates the speedups and peak intermediate memory
of every app. A second runs an ablation of the CPU apps that adds the
optimizations inductive functions enable one at a time. Every form is checked
against a reference implementation before it is timed.

\subsection{Artifact check-list (meta-information)}

{\small
\begin{itemize}
  \item {\bf Algorithm:} inductive functions in Halide, scheduled with fusion,
    sliding window and storage folding.
  \item {\bf Program:} six CPU apps (RNG, biquads, sequence alignment, prefix
    mean, Viterbi, ODE) and two GPU apps (FlashAttention, Mamba-2 forward
    and backward), with their third-party baselines.
  \item {\bf Compilation:} Halide built against LLVM~22, which \texttt{uv}
    installs; C++17 at \texttt{-O3 -march=native}; CUDA~13 for the GPU apps.
  \item {\bf Data set:} generated by each app from a fixed seed.
  \item {\bf Run-time environment:} Linux x86-64 with jemalloc. The optional
    baselines are Boost.odeint, oneTBB, Intel IPP, parasail, Julia, scipy,
    PyTorch, Triton and mamba\_ssm.
  \item {\bf Hardware:} a multicore x86-64 CPU. The paper's machine
    (\artsec{sec:benchmarks}{in the paper}) has 32 cores with AVX-512; AVX2 also
    works. The GPU apps need an NVIDIA GPU with tensor cores.
  \item {\bf Execution:} two shell scripts. All-cores rows use one thread
    per physical core.
  \item {\bf Metrics:} best of 30 timed runs after 3 warmups (ms); peak
    intermediate heap (MB) from the Halide profiler.
  \item {\bf Output:} Markdown and text tables, plus logs.
  \item {\bf Experiments:} \artfig{fig:cpuspeedup}{CPU speedups},
    \artfig{fig:gpuspeedup}{GPU speedups}, \artfig{fig:memuse}{memory}, and
    the CPU ablation.
  \item {\bf How much disk space required (approximately)?:} 10~GB.
  \item {\bf How much time is needed to prepare workflow (approximately)?:}
    about an hour.
  \item {\bf How much time is needed to complete experiments (approximately)?:}
    a few hours; each app can also be run on its own.
  \item {\bf Publicly available?:} yes.
  \item {\bf Code licenses (if publicly available)?:} MIT (Halide). Vendored
    baselines keep their own licenses.
  \item {\bf Archived (provide DOI)?:} \artifactdoi.
\end{itemize}
}

\subsection{Description}

\subsubsection{How delivered}
The branch \texttt{inductive\_artifact} of \artifactrepo, archived at
\artifactdoi. The benchmark apps are in \path{apps/}, and the artifact's
scripts and instructions are in \path{inductive_artifact/}.

\subsubsection{Hardware dependencies}
The CPU apps run on any x86-64 Linux machine. The all-cores rows are
memory-bound, so their speedups depend on the memory bandwidth available per
core. The GPU apps need an NVIDIA GPU with tensor cores and CUDA~13.

\subsubsection{Software dependencies}
Linux, a C++17 compiler, GNU make, pkg-config, \texttt{uv}, jemalloc, and the
libpng, libjpeg and zlib headers. Every process the scripts launch, baselines
included, preloads jemalloc configured to keep freed memory, as in the paper.
The baselines are optional. Boost.odeint, oneTBB, Intel IPP and parasail are
found through \texttt{BOOST\_DIR}, \texttt{TBB\_DIR}, \texttt{IPP\_DIR} and
\texttt{PARASAIL\_DIR}, or a system install. Julia is found on the
\texttt{PATH}. The Python baselines come from
\path{inductive_artifact/requirements.txt}. When a baseline is missing, the
row reports the fastest baseline that is present.

\subsubsection{Data sets}
None; every input is generated.

\subsection{Installation}

{\footnotesize
\begin{verbatim}
git clone -b inductive_artifact \
  https://github.com/wraith1995/Halide.git
cd Halide
sudo apt install build-essential pkg-config \
  libpng-dev libjpeg-turbo8-dev zlib1g-dev \
  libjemalloc2 libboost-dev libtbb-dev
inductive_artifact/setup.sh
# or, with the Python and GPU baselines:
WITH_PYTHON_BASELINES=1 inductive_artifact/setup.sh
\end{verbatim}
}

\noindent\texttt{setup.sh} installs LLVM into \path{.venv} and builds the
Halide distribution, \path{distrib/}, which every app links against.

\subsection{Experiment workflow}

{\footnotesize
\begin{verbatim}
inductive_artifact/run_benchmarks.sh  # figures
inductive_artifact/run_ablation.sh    # ablation
\end{verbatim}
}

\noindent Both scripts accept \texttt{--only} with a comma-separated list of
apps: \texttt{cpu\_rng}, \texttt{cpu\_biquads}, \texttt{cpu\_alignment},
\texttt{prefixsum}, \texttt{viterbi}, \texttt{ode}, \texttt{flash\_attention}
and \texttt{cuda\_mamba2}. Finished rows are cached, so a partial run
updates the table rather than replacing it. The GPU apps print their own
ablation levels:

{\footnotesize
\begin{verbatim}
make -C apps/cuda_mamba2 SCAN=inductive \
  WMMA=true ABLATE_FUSION=1 test test_bwd
make -C apps/cuda_attention test
\end{verbatim}
}

\subsection{Evaluation and expected results}

\path{run_benchmarks.sh} writes \path{inductive_artifact/results/benchmarks.md},
with one row per app and configuration: the inductive and non-inductive
times, each form's peak intermediate memory, and the fastest baseline. The
all-cores CPU rows and the GPU rows are the configurations in the paper. The
single-thread rows and the Chebyshev row are controls that the paper does
not plot. The tables from the paper's machine are in
\path{apps/inductive_benchmarks_results.md} and
\path{apps/inductive_ablation_results/}. The claims to check:

\begin{enumerate}
  \item \textbf{Correctness.} Every form passes its app's check against the
    reference before it is timed.
  \item \textbf{Memory (\artfig{fig:memuse}{memory figure}).} The inductive
    forms allocate no intermediate memory, except for the backpointers that
    alignment (1094~MB) and Viterbi (134~MB) keep for their traceback. The
    non-inductive forms materialize the trajectory: 4295~MB for RNG,
    2147~MB for biquads, 10758~MB for alignment, 1074~MB each for prefix
    mean and ODE, and 671~MB for Viterbi. On the GPU, Mamba-2 forward needs
    42~MB against 141~MB. These values depend only on the problem sizes, so
    they should match exactly. The ODE and Viterbi all-cores rows run one
    batch element per core, so their memory scales with the core count;
    \texttt{HB\_NCORES=32} gives the paper's sizes.
  \item \textbf{CPU speedups (\artfig{fig:cpuspeedup}{CPU figure}).} The
    inductive form is faster than the non-inductive form in every app. On
    the paper's machine the speedups are prefix mean 2.00$\times$, Viterbi
    4.37$\times$, RNG 5.88$\times$, biquads 7.43$\times$, alignment
    16.47$\times$ and ODE 19.01$\times$. Against the baselines, the inductive
    form matches Julia (RNG) and oneTBB (prefix mean), reaches 93\% of
    Finding Fast Filters (biquads), and beats parasail (alignment) and
    Boost.odeint (ODE).
  \item \textbf{GPU speedups (\artfig{fig:gpuspeedup}{GPU figure}).} The
    inductive form beats the non-inductive form by 1.09$\times$
    (FlashAttention), 1.15$\times$ (Mamba-2 forward) and 1.04$\times$
    (Mamba-2 backward). Relative to the vendor kernels it reaches 94\% of
    PyTorch, 98\% of Triton on the forward pass and 73\% of Triton on the
    backward pass.
  \item \textbf{Ablation (CPU).} The levels are non-inductive, fusion only,
    fusion with a sliding window, and the full schedule with storage
    folding. Fusion alone is far slower than any other level, because each
    output recomputes the recurrence from its start. It is measured at a
    small size, where it runs 62--6038$\times$ slower than the full schedule.
    Adding the sliding window lands near the non-inductive form
    (1.3--14$\times$ the full schedule). ODE is the exception at
    39$\times$: its unfolded window is three times the non-inductive
    storage. Storage folding is the fastest level in every app.
\end{enumerate}

Absolute times depend on the machine. The all-cores speedups come from
memory traffic, so they are smaller on a machine with more memory bandwidth
per core. The ranking of the apps and the memory values should not change.

\subsection{Experiment customization}

Each app's Makefile takes its problem size on the command line, and the
README in each app lists the knobs. Environment variables set the number of
cores (\texttt{HB\_NCORES}), the core that single-thread rows are pinned to
(\texttt{HB\_PIN\_CORE}), the Python used for the baselines
(\texttt{HB\_PYTHON}), and the jemalloc library (\texttt{HB\_JEMALLOC}).

\subsection{Methodology}

Submission, reviewing and badging methodology:
\begin{itemize}
  \item \url{https://www.acm.org/publications/policies/artifact-review-and-badging-current}
  \item \url{https://cTuning.org/ae}
\end{itemize}
 before the bibliography (CGO
% places the artifact appendix before the References, up to two pages).
% After archiving the artifact, set its DOI here:
\newcommand{\artifactdoi}{\textit{DOI to be assigned on archiving}}
\newcommand{\artifactrepo}{\url{https://github.com/wraith1995/Halide/tree/inductive_artifact}}
% Figure references that fall back to text when built outside the paper.
\newcommand{\artfig}[2]{\ifcsname r@#1\endcsname Figure~\ref{#1}\else #2\fi}
\newcommand{\artsec}[2]{\ifcsname r@#1\endcsname Section~\ref{#1}\else #2\fi}

\section{Artifact Appendix}
\sloppy

\subsection{Abstract}

The artifact is a fork of Halide with inductive functions, and the benchmark
suite behind \artfig{fig:cpuspeedup}{the CPU speedup figure},
\artfig{fig:gpuspeedup}{the GPU speedup figure} and
\artfig{fig:memuse}{the memory figure}. Each benchmark computes the same
recurrence twice: as an inductive function, and as a non-inductive
update definition over an RDom, which materializes the whole trajectory.
Each run also times the fastest third-party implementation that is
installed. One script regenerates the speedups and peak intermediate memory
of every app. A second runs an ablation of the CPU apps that adds the
optimizations inductive functions enable one at a time. Every form is checked
against a reference implementation before it is timed.

\subsection{Artifact check-list (meta-information)}

{\small
\begin{itemize}
  \item {\bf Algorithm:} inductive functions in Halide, scheduled with fusion,
    sliding window and storage folding.
  \item {\bf Program:} six CPU apps (RNG, biquads, sequence alignment, prefix
    mean, Viterbi, ODE) and two GPU apps (FlashAttention, Mamba-2 forward
    and backward), with their third-party baselines.
  \item {\bf Compilation:} Halide built against LLVM~22, which \texttt{uv}
    installs; C++17 at \texttt{-O3 -march=native}; CUDA~13 for the GPU apps.
  \item {\bf Data set:} generated by each app from a fixed seed.
  \item {\bf Run-time environment:} Linux x86-64 with jemalloc. The optional
    baselines are Boost.odeint, oneTBB, Intel IPP, parasail, Julia, scipy,
    PyTorch, Triton and mamba\_ssm.
  \item {\bf Hardware:} a multicore x86-64 CPU. The paper's machine
    (\artsec{sec:benchmarks}{in the paper}) has 32 cores with AVX-512; AVX2 also
    works. The GPU apps need an NVIDIA GPU with tensor cores.
  \item {\bf Execution:} two shell scripts. All-cores rows use one thread
    per physical core.
  \item {\bf Metrics:} best of 30 timed runs after 3 warmups (ms); peak
    intermediate heap (MB) from the Halide profiler.
  \item {\bf Output:} Markdown and text tables, plus logs.
  \item {\bf Experiments:} \artfig{fig:cpuspeedup}{CPU speedups},
    \artfig{fig:gpuspeedup}{GPU speedups}, \artfig{fig:memuse}{memory}, and
    the CPU ablation.
  \item {\bf How much disk space required (approximately)?:} 10~GB.
  \item {\bf How much time is needed to prepare workflow (approximately)?:}
    about an hour.
  \item {\bf How much time is needed to complete experiments (approximately)?:}
    a few hours; each app can also be run on its own.
  \item {\bf Publicly available?:} yes.
  \item {\bf Code licenses (if publicly available)?:} MIT (Halide). Vendored
    baselines keep their own licenses.
  \item {\bf Archived (provide DOI)?:} \artifactdoi.
\end{itemize}
}

\subsection{Description}

\subsubsection{How delivered}
The branch \texttt{inductive\_artifact} of \artifactrepo, archived at
\artifactdoi. The benchmark apps are in \path{apps/}, and the artifact's
scripts and instructions are in \path{inductive_artifact/}.

\subsubsection{Hardware dependencies}
The CPU apps run on any x86-64 Linux machine. The all-cores rows are
memory-bound, so their speedups depend on the memory bandwidth available per
core. The GPU apps need an NVIDIA GPU with tensor cores and CUDA~13.

\subsubsection{Software dependencies}
Linux, a C++17 compiler, GNU make, pkg-config, \texttt{uv}, jemalloc, and the
libpng, libjpeg and zlib headers. Every process the scripts launch, baselines
included, preloads jemalloc configured to keep freed memory, as in the paper.
The baselines are optional. Boost.odeint, oneTBB, Intel IPP and parasail are
found through \texttt{BOOST\_DIR}, \texttt{TBB\_DIR}, \texttt{IPP\_DIR} and
\texttt{PARASAIL\_DIR}, or a system install. Julia is found on the
\texttt{PATH}. The Python baselines come from
\path{inductive_artifact/requirements.txt}. When a baseline is missing, the
row reports the fastest baseline that is present.

\subsubsection{Data sets}
None; every input is generated.

\subsection{Installation}

{\footnotesize
\begin{verbatim}
git clone -b inductive_artifact \
  https://github.com/wraith1995/Halide.git
cd Halide
sudo apt install build-essential pkg-config \
  libpng-dev libjpeg-turbo8-dev zlib1g-dev \
  libjemalloc2 libboost-dev libtbb-dev
inductive_artifact/setup.sh
# or, with the Python and GPU baselines:
WITH_PYTHON_BASELINES=1 inductive_artifact/setup.sh
\end{verbatim}
}

\noindent\texttt{setup.sh} installs LLVM into \path{.venv} and builds the
Halide distribution, \path{distrib/}, which every app links against.

\subsection{Experiment workflow}

{\footnotesize
\begin{verbatim}
inductive_artifact/run_benchmarks.sh  # figures
inductive_artifact/run_ablation.sh    # ablation
\end{verbatim}
}

\noindent Both scripts accept \texttt{--only} with a comma-separated list of
apps: \texttt{cpu\_rng}, \texttt{cpu\_biquads}, \texttt{cpu\_alignment},
\texttt{prefixsum}, \texttt{viterbi}, \texttt{ode}, \texttt{flash\_attention}
and \texttt{cuda\_mamba2}. Finished rows are cached, so a partial run
updates the table rather than replacing it. The GPU apps print their own
ablation levels:

{\footnotesize
\begin{verbatim}
make -C apps/cuda_mamba2 SCAN=inductive \
  WMMA=true ABLATE_FUSION=1 test test_bwd
make -C apps/cuda_attention test
\end{verbatim}
}

\subsection{Evaluation and expected results}

\path{run_benchmarks.sh} writes \path{inductive_artifact/results/benchmarks.md},
with one row per app and configuration: the inductive and non-inductive
times, each form's peak intermediate memory, and the fastest baseline. The
all-cores CPU rows and the GPU rows are the configurations in the paper. The
single-thread rows and the Chebyshev row are controls that the paper does
not plot. The tables from the paper's machine are in
\path{apps/inductive_benchmarks_results.md} and
\path{apps/inductive_ablation_results/}. The claims to check:

\begin{enumerate}
  \item \textbf{Correctness.} Every form passes its app's check against the
    reference before it is timed.
  \item \textbf{Memory (\artfig{fig:memuse}{memory figure}).} The inductive
    forms allocate no intermediate memory, except for the backpointers that
    alignment (1094~MB) and Viterbi (134~MB) keep for their traceback. The
    non-inductive forms materialize the trajectory: 4295~MB for RNG,
    2147~MB for biquads, 10758~MB for alignment, 1074~MB each for prefix
    mean and ODE, and 671~MB for Viterbi. On the GPU, Mamba-2 forward needs
    42~MB against 141~MB. These values depend only on the problem sizes, so
    they should match exactly. The ODE and Viterbi all-cores rows run one
    batch element per core, so their memory scales with the core count;
    \texttt{HB\_NCORES=32} gives the paper's sizes.
  \item \textbf{CPU speedups (\artfig{fig:cpuspeedup}{CPU figure}).} The
    inductive form is faster than the non-inductive form in every app. On
    the paper's machine the speedups are prefix mean 2.00$\times$, Viterbi
    4.37$\times$, RNG 5.88$\times$, biquads 7.43$\times$, alignment
    16.47$\times$ and ODE 19.01$\times$. Against the baselines, the inductive
    form matches Julia (RNG) and oneTBB (prefix mean), reaches 93\% of
    Finding Fast Filters (biquads), and beats parasail (alignment) and
    Boost.odeint (ODE).
  \item \textbf{GPU speedups (\artfig{fig:gpuspeedup}{GPU figure}).} The
    inductive form beats the non-inductive form by 1.09$\times$
    (FlashAttention), 1.15$\times$ (Mamba-2 forward) and 1.04$\times$
    (Mamba-2 backward). Relative to the vendor kernels it reaches 94\% of
    PyTorch, 98\% of Triton on the forward pass and 73\% of Triton on the
    backward pass.
  \item \textbf{Ablation (CPU).} The levels are non-inductive, fusion only,
    fusion with a sliding window, and the full schedule with storage
    folding. Fusion alone is far slower than any other level, because each
    output recomputes the recurrence from its start. It is measured at a
    small size, where it runs 62--6038$\times$ slower than the full schedule.
    Adding the sliding window lands near the non-inductive form
    (1.3--14$\times$ the full schedule). ODE is the exception at
    39$\times$: its unfolded window is three times the non-inductive
    storage. Storage folding is the fastest level in every app.
\end{enumerate}

Absolute times depend on the machine. The all-cores speedups come from
memory traffic, so they are smaller on a machine with more memory bandwidth
per core. The ranking of the apps and the memory values should not change.

\subsection{Experiment customization}

Each app's Makefile takes its problem size on the command line, and the
README in each app lists the knobs. Environment variables set the number of
cores (\texttt{HB\_NCORES}), the core that single-thread rows are pinned to
(\texttt{HB\_PIN\_CORE}), the Python used for the baselines
(\texttt{HB\_PYTHON}), and the jemalloc library (\texttt{HB\_JEMALLOC}).

\subsection{Methodology}

Submission, reviewing and badging methodology:
\begin{itemize}
  \item \url{https://www.acm.org/publications/policies/artifact-review-and-badging-current}
  \item \url{https://cTuning.org/ae}
\end{itemize}
 before the bibliography (CGO
% places the artifact appendix before the References, up to two pages).
% After archiving the artifact, set its DOI here:
\newcommand{\artifactdoi}{\textit{DOI to be assigned on archiving}}
\newcommand{\artifactrepo}{\url{https://github.com/wraith1995/Halide/tree/inductive_artifact}}
% Figure references that fall back to text when built outside the paper.
\newcommand{\artfig}[2]{\ifcsname r@#1\endcsname Figure~\ref{#1}\else #2\fi}
\newcommand{\artsec}[2]{\ifcsname r@#1\endcsname Section~\ref{#1}\else #2\fi}

\section{Artifact Appendix}
\sloppy

\subsection{Abstract}

The artifact is a fork of Halide with inductive functions, and the benchmark
suite behind \artfig{fig:cpuspeedup}{the CPU speedup figure},
\artfig{fig:gpuspeedup}{the GPU speedup figure} and
\artfig{fig:memuse}{the memory figure}. Each benchmark computes the same
recurrence twice: as an inductive function, and as a non-inductive
update definition over an RDom, which materializes the whole trajectory.
Each run also times the fastest third-party implementation that is
installed. One script regenerates the speedups and peak intermediate memory
of every app. A second runs an ablation of the CPU apps that adds the
optimizations inductive functions enable one at a time. Every form is checked
against a reference implementation before it is timed.

\subsection{Artifact check-list (meta-information)}

{\small
\begin{itemize}
  \item {\bf Algorithm:} inductive functions in Halide, scheduled with fusion,
    sliding window and storage folding.
  \item {\bf Program:} six CPU apps (RNG, biquads, sequence alignment, prefix
    mean, Viterbi, ODE) and two GPU apps (FlashAttention, Mamba-2 forward
    and backward), with their third-party baselines.
  \item {\bf Compilation:} Halide built against LLVM~22, which \texttt{uv}
    installs; C++17 at \texttt{-O3 -march=native}; CUDA~13 for the GPU apps.
  \item {\bf Data set:} generated by each app from a fixed seed.
  \item {\bf Run-time environment:} Linux x86-64 with jemalloc. The optional
    baselines are Boost.odeint, oneTBB, Intel IPP, parasail, Julia, scipy,
    PyTorch, Triton and mamba\_ssm.
  \item {\bf Hardware:} a multicore x86-64 CPU. The paper's machine
    (\artsec{sec:benchmarks}{in the paper}) has 32 cores with AVX-512; AVX2 also
    works. The GPU apps need an NVIDIA GPU with tensor cores.
  \item {\bf Execution:} two shell scripts. All-cores rows use one thread
    per physical core.
  \item {\bf Metrics:} best of 30 timed runs after 3 warmups (ms); peak
    intermediate heap (MB) from the Halide profiler.
  \item {\bf Output:} Markdown and text tables, plus logs.
  \item {\bf Experiments:} \artfig{fig:cpuspeedup}{CPU speedups},
    \artfig{fig:gpuspeedup}{GPU speedups}, \artfig{fig:memuse}{memory}, and
    the CPU ablation.
  \item {\bf How much disk space required (approximately)?:} 10~GB.
  \item {\bf How much time is needed to prepare workflow (approximately)?:}
    about an hour.
  \item {\bf How much time is needed to complete experiments (approximately)?:}
    a few hours; each app can also be run on its own.
  \item {\bf Publicly available?:} yes.
  \item {\bf Code licenses (if publicly available)?:} MIT (Halide). Vendored
    baselines keep their own licenses.
  \item {\bf Archived (provide DOI)?:} \artifactdoi.
\end{itemize}
}

\subsection{Description}

\subsubsection{How delivered}
The branch \texttt{inductive\_artifact} of \artifactrepo, archived at
\artifactdoi. The benchmark apps are in \path{apps/}, and the artifact's
scripts and instructions are in \path{inductive_artifact/}.

\subsubsection{Hardware dependencies}
The CPU apps run on any x86-64 Linux machine. The all-cores rows are
memory-bound, so their speedups depend on the memory bandwidth available per
core. The GPU apps need an NVIDIA GPU with tensor cores and CUDA~13.

\subsubsection{Software dependencies}
Linux, a C++17 compiler, GNU make, pkg-config, \texttt{uv}, jemalloc, and the
libpng, libjpeg and zlib headers. Every process the scripts launch, baselines
included, preloads jemalloc configured to keep freed memory, as in the paper.
The baselines are optional. Boost.odeint, oneTBB, Intel IPP and parasail are
found through \texttt{BOOST\_DIR}, \texttt{TBB\_DIR}, \texttt{IPP\_DIR} and
\texttt{PARASAIL\_DIR}, or a system install. Julia is found on the
\texttt{PATH}. The Python baselines come from
\path{inductive_artifact/requirements.txt}. When a baseline is missing, the
row reports the fastest baseline that is present.

\subsubsection{Data sets}
None; every input is generated.

\subsection{Installation}

{\footnotesize
\begin{verbatim}
git clone -b inductive_artifact \
  https://github.com/wraith1995/Halide.git
cd Halide
sudo apt install build-essential pkg-config \
  libpng-dev libjpeg-turbo8-dev zlib1g-dev \
  libjemalloc2 libboost-dev libtbb-dev
inductive_artifact/setup.sh
# or, with the Python and GPU baselines:
WITH_PYTHON_BASELINES=1 inductive_artifact/setup.sh
\end{verbatim}
}

\noindent\texttt{setup.sh} installs LLVM into \path{.venv} and builds the
Halide distribution, \path{distrib/}, which every app links against.

\subsection{Experiment workflow}

{\footnotesize
\begin{verbatim}
inductive_artifact/run_benchmarks.sh  # figures
inductive_artifact/run_ablation.sh    # ablation
\end{verbatim}
}

\noindent Both scripts accept \texttt{--only} with a comma-separated list of
apps: \texttt{cpu\_rng}, \texttt{cpu\_biquads}, \texttt{cpu\_alignment},
\texttt{prefixsum}, \texttt{viterbi}, \texttt{ode}, \texttt{flash\_attention}
and \texttt{cuda\_mamba2}. Finished rows are cached, so a partial run
updates the table rather than replacing it. The GPU apps print their own
ablation levels:

{\footnotesize
\begin{verbatim}
make -C apps/cuda_mamba2 SCAN=inductive \
  WMMA=true ABLATE_FUSION=1 test test_bwd
make -C apps/cuda_attention test
\end{verbatim}
}

\subsection{Evaluation and expected results}

\path{run_benchmarks.sh} writes \path{inductive_artifact/results/benchmarks.md},
with one row per app and configuration: the inductive and non-inductive
times, each form's peak intermediate memory, and the fastest baseline. The
all-cores CPU rows and the GPU rows are the configurations in the paper. The
single-thread rows and the Chebyshev row are controls that the paper does
not plot. The tables from the paper's machine are in
\path{apps/inductive_benchmarks_results.md} and
\path{apps/inductive_ablation_results/}. The claims to check:

\begin{enumerate}
  \item \textbf{Correctness.} Every form passes its app's check against the
    reference before it is timed.
  \item \textbf{Memory (\artfig{fig:memuse}{memory figure}).} The inductive
    forms allocate no intermediate memory, except for the backpointers that
    alignment (1094~MB) and Viterbi (134~MB) keep for their traceback. The
    non-inductive forms materialize the trajectory: 4295~MB for RNG,
    2147~MB for biquads, 10758~MB for alignment, 1074~MB each for prefix
    mean and ODE, and 671~MB for Viterbi. On the GPU, Mamba-2 forward needs
    42~MB against 141~MB. These values depend only on the problem sizes, so
    they should match exactly. The ODE and Viterbi all-cores rows run one
    batch element per core, so their memory scales with the core count;
    \texttt{HB\_NCORES=32} gives the paper's sizes.
  \item \textbf{CPU speedups (\artfig{fig:cpuspeedup}{CPU figure}).} The
    inductive form is faster than the non-inductive form in every app. On
    the paper's machine the speedups are prefix mean 2.00$\times$, Viterbi
    4.37$\times$, RNG 5.88$\times$, biquads 7.43$\times$, alignment
    16.47$\times$ and ODE 19.01$\times$. Against the baselines, the inductive
    form matches Julia (RNG) and oneTBB (prefix mean), reaches 93\% of
    Finding Fast Filters (biquads), and beats parasail (alignment) and
    Boost.odeint (ODE).
  \item \textbf{GPU speedups (\artfig{fig:gpuspeedup}{GPU figure}).} The
    inductive form beats the non-inductive form by 1.09$\times$
    (FlashAttention), 1.15$\times$ (Mamba-2 forward) and 1.04$\times$
    (Mamba-2 backward). Relative to the vendor kernels it reaches 94\% of
    PyTorch, 98\% of Triton on the forward pass and 73\% of Triton on the
    backward pass.
  \item \textbf{Ablation (CPU).} The levels are non-inductive, fusion only,
    fusion with a sliding window, and the full schedule with storage
    folding. Fusion alone is far slower than any other level, because each
    output recomputes the recurrence from its start. It is measured at a
    small size, where it runs 62--6038$\times$ slower than the full schedule.
    Adding the sliding window lands near the non-inductive form
    (1.3--14$\times$ the full schedule). ODE is the exception at
    39$\times$: its unfolded window is three times the non-inductive
    storage. Storage folding is the fastest level in every app.
\end{enumerate}

Absolute times depend on the machine. The all-cores speedups come from
memory traffic, so they are smaller on a machine with more memory bandwidth
per core. The ranking of the apps and the memory values should not change.

\subsection{Experiment customization}

Each app's Makefile takes its problem size on the command line, and the
README in each app lists the knobs. Environment variables set the number of
cores (\texttt{HB\_NCORES}), the core that single-thread rows are pinned to
(\texttt{HB\_PIN\_CORE}), the Python used for the baselines
(\texttt{HB\_PYTHON}), and the jemalloc library (\texttt{HB\_JEMALLOC}).

\subsection{Methodology}

Submission, reviewing and badging methodology:
\begin{itemize}
  \item \url{https://www.acm.org/publications/policies/artifact-review-and-badging-current}
  \item \url{https://cTuning.org/ae}
\end{itemize}
